\documentclass[10pt,a4paper]{amsart}
\usepackage[margin=19mm]{geometry}
\usepackage{amsmath,amssymb,mathtools}
\usepackage[scr=rsfs,cal=euler]{mathalfa}
\usepackage{xfrac}
\usepackage[unicode,hidelinks,bookmarks=false]{hyperref}
\newcommand{\ie}{i.e.\ }
\newcommand{\cf}{cf.\ }
\newcommand{\resp}{respectively\ }
\newcommand{\Subsection}{\subsection}
\providecommand{\nobreakdash}{\nobreak\hbox{-}}
\setlength{\emergencystretch}{2em}
\multlinegap0pt
\usepackage{bm}
\usepackage{bbm}
\usepackage{dsfont}
\usepackage{xspace}
\usepackage{cancel}
\usepackage{mathtools}
\usepackage{amsthm}
\makeatletter
\@ifundefined{theorem}{\newtheorem{theorem}{Theorem}}{}
\makeatother
\newcommand{\modulo}[1]{\textup{ (mod }#1)}
\newcommand{\quot}[2]{\large\sfrac{#1}{#2}\normalsize}
\title[Multiplicative Relationships of Subrings and their Applicati]{Multiplicative Relationships of Subrings and their Applications to Factorization}
\author{Grant Moles}
\date{Source: arXiv:2506.24031v2 (CC BY 4.0)}
\begin{document}
\maketitle

\noindent\emph{Source and licence.} Multiplicative Relationships of Subrings and their Applications to Factorization. Version arXiv:2506.24031v2 (CC BY 4.0), distributed under the Creative Commons Attribution 4.0 International licence, as recorded in the arXiv licence metadata for this exact version and shown on the abstract page https://arxiv.org/abs/2506.24031v2. Full rights notice: licenses/arxiv-2506.24031-CC-BY-4.0.txt.\par\smallskip

\noindent\textit{Editorial scope.} Counted unit: the Theorem whose source label is \texttt{ip equiv conditions} in the article (source0.tex lines 320--343, source label \texttt{ip equiv conditions}), delivered together with the article's own complete original proof of it (source0.tex lines 345--365, one \texttt{proof} environment of 21 lines). No mathematics is written, repaired or re-derived: statement and proof are the article's own text line for line. Required context retained uncounted: none. Every definition, display and label of those lines is retained unchanged. Omitted from this package and not redistributed: 1--319, 344--344, 366--764. Nothing inside a retained excerpt is omitted or altered: every retained body is delivered exactly as the article's lines are written, so no omission receipt of any kind accompanies this package. Citations: the retained excerpts carry no citation command at all, so no bibliography entry of the article is transcribed and no reference is redistributed. Reference adaptation: none was needed and none was made; every reference inside the delivered unit resolves inside it, so no unresolved reference or citation is produced. Printed slips kept literal: no printed word, symbol, hypothesis or number of the counted statement or of its proof was altered. Layout and class: the article's own class is \texttt{amsart} with packages including geometry, url, amsmath, amssymb, mathrsfs, enumitem, commath, physics, nicefrac, amsthm, xfrac, xy, graphicx; the wrapper is the lane's pinned \texttt{amsart} editorial wrapper at 10pt on A4, which restates the theorem-like environments the retained text needs and adds the article's own macro definitions that the retained text needs (\texttt{\textbackslash modulo}, \texttt{\textbackslash quot}), with extra wrapper packages (bm, bbm, dsfont, xspace, cancel, mathtools). The delivered numbering is the unit's own. Assets: no retained span contains a figure environment, a graphics-inclusion command, a PDF-page inclusion, a table or a TikZ picture; no image file is copied into this package, the package declares no asset dependency and no image file is redistributed with it. Licence: version arXiv:2506.24031v2 of this article is distributed under the Creative Commons Attribution 4.0 International licence, as recorded in the arXiv licence metadata for this exact version and shown on the abstract page https://arxiv.org/abs/2506.24031v2. A full rights notice is delivered at licenses/arxiv-2506.24031-CC-BY-4.0.txt. Characters: every retained excerpt is pure ASCII, so the delivered engine, XeLaTeX with its default Latin Modern text font, reports no missing character.
	\begin{theorem}
		\label{ip equiv conditions}
		Let $T$ be a commutative ring, $R\subseteq T$ a subring (neither assumed to have identity), and $I=(R:T)$. Consider the following condition:
		\begin{enumerate}
			\item $R$ is an ideal-preserving subring of $T$.
		\end{enumerate}
		If $T$ is an integral domain, the following is equivalent to (1):
		\begin{enumerate}
			\item[(2)] For any $T$-ideal $J$ and $t\in T$, $$t\notin J\implies R\cap (t)\nsubseteq J.$$
		\end{enumerate}
		If $T$ is a Dedekind domain, the following is equivalent to (1):
		\begin{enumerate}
			\item[(3)] If $P_1,P_2$ are prime $T$-ideals, then $R\cap P_1\nsubseteq P_1^2$ and, if $P_1\neq P_2$, then $R\cap P_1\nsubseteq P_2$.
		\end{enumerate}
		If $T$ is a Dedekind domain and $1\in R$, the following is equivalent to (1):
		\begin{enumerate}
			\item[(4)] If $P_1,P_2$ are prime $T$-ideals containing $I$, then $R\cap P_1\nsubseteq P_1^2$, and if $P_1\neq P_2$, then $R\cap P_1\nsubseteq P_2$.
		\end{enumerate}
		If $T$ is a Dedekind domain, $1\in R$, and $T$ is integral over $R$, then the following is equivalent to (1):
		\begin{enumerate}
			\item[(5)] If $I=P_1^{a_1}\dots P_k^{a_k}$ is the factorization of $I$ into prime $T$-ideals, then $R\cap P_i\nsubseteq P_i^2$ for $1\leq i\leq k$ and 
			$$\quot{R}{I}\cong \prod_{i=1}^k\quot{R}{R\cap P_i^{a_i}}\cong \prod_{i=1}^k\quot{R+P_i^{a_i}}{P_i^{a_i}}.$$
		\end{enumerate}
	\end{theorem}

	\begin{proof}
		Suppose that $T$ is an integral domain. In this case, we need to show $1\iff 2$. Assume that Condition 1 holds and let $t\in T$ and $J$ be a $T$-ideal. Then note that $t\notin J$ is equivalent to the statement $(t)\nsubseteq J$. Since $R$ is an ideal-preserving subring, this would imply that $R\cap (t)\nsubseteq J$. Then $1\implies 2$.
		
		Now assume that Condition 2 holds and let $J_1, J_2$ be $T$-ideals such that $J_1\nsubseteq J_2$. Then there is some $t\in J_1\backslash J_2$. Since $t\notin J_2$, Condition 2 tells us that $R\cap (t)\nsubseteq J_2$. Then since $R\cap (t)\subseteq R\cap J_1$, it follows that $R\cap J_1\nsubseteq J_2$. Then $R$ is an ideal-preserving subring of $T$, so $2\implies 1$.
		
		Now suppose that $T$ is a Dedekind domain; we want to show that $1\iff 3$. It is immediately clear in this case that $1\implies 3$, since the statement of $3$ is simply an application of the definition of an ideal-preserving subring to the prime ideals in $T$. We need to show the converse. To that end, let $J_1\nsubseteq J_2$ be $T$-ideals. Since $T$ is a Dedekind domain, this is equivalent to the statement $J_2\nmid J_1$. Then letting $J_1=P_1^{a_1}\dots P_k^{a_k}$ and $J_2=P_1^{b_1}\dots P_k^{b_k}$ be the representations of $J_1$ and $J_2$ as products of prime $T$-ideals (allowing the possibility that some $a_i$ or $b_j$ is 0), it must be the case that for some $1\leq i\leq k$, $b_i>a_i$. By Condition 3, we can select some element $\alpha_i\in R\cap P_i\backslash P_i^2$ and, for each $1\leq j\leq k$, $j\neq i$, we can select $\alpha_j\in R\cap P_j\backslash P_i$. Then the element $\alpha=\alpha_1^{a_1}\dots \alpha_k^{a_k}\in R\cap (P_1^{a_1}\dots P_k^{a_k})=R\cap J_1$. However, by the choice of the $\alpha_j$'s, $P_i^{a_i}$ is the exact power of $P_i$ dividing $(\alpha)$. Then since $b_i>a_i$, $\alpha\notin P_i^{b_i}\supseteq J_2$, so $\alpha\in R\cap J_1\backslash J_2$. Therefore, $R$ is an ideal-preserving subring of $T$, so $3\implies 1$.
		
		Next, suppose that $T$ is a Dedekind domain and $1\in R$. We want to show $1\iff 4$; since $T$ is a Dedekind domain, it will suffice to show $3\iff 4$. Again, $3\implies 4$ trivially, since we are only limiting the prime $T$-ideals we are considering. We need only show the converse, $4\implies 3$. To do so, we will show that Condition 3 will always hold if either $P_1$ or $P_2$ (or both) do not divide the conductor ideal. It is worth noting at this point that we will be using the convention that any ideal divides (contains) the zero ideal; thus, if $I=\{0\}$, Conditions 3 and 4 are exactly the same.
		
		First, suppose that $P_1$ is a prime $T$-ideal such that $P_1\nmid I$ (if $I$ has a prime divisor; if not, then $T=R=I$ and $R$ is trivially an ideal-preserving subring of $T$). Then letting $\alpha\in P_1\backslash P_1^2$ and $\beta\in I\backslash P_1$, we have $\alpha\beta\in R\cap P_1\backslash P_1^2$, so $R\cap P_1\nsubseteq P_1^2$. Now let $P_1$ and $P_2$ be distinct prime $T$-ideals. If $P_2\nmid I$, then let $\alpha\in P_1\backslash P_2$ and $\beta\in I\backslash P_2$. Then $\alpha\beta\in R\cap P_1\backslash P_2$, so $R\cap P_1\nsubseteq P_2$. 
		
		The only remaining case to consider is when $P_1\nmid I$ and $P_2|I$. In this case, note that $P_1$ and $I$ are relatively prime ideals in $T$, and thus there exist $\alpha\in P_1$ and $\beta\in I$ such that $1=\alpha+\beta$. Then note that $\alpha=1-\beta\in P_1$. Since $\beta\in I\subseteq P_2$, $\alpha=1-\beta\notin P_2$. Finally, since $1\in R$ and $\beta\in I\subseteq R$, $\alpha=1-\beta\in R$. Then $\alpha=1-\beta\in R\cap P_1\backslash P_2$, so $R\cap P_1\nsubseteq P_2$. Then $4\implies 3$.
		
		Finally, assume that $T$ is a Dedekind domain, $1\in R$, and $T$ is integral over $R$. We want to show $1\iff 5$. To start, we will note that the second isomorphism in the statement of Condition 5 will always hold, regardless of whether $R$ is an ideal-preserving subring of $T$. This can easily be seen by considering the obvious mapping $\phi:\prod_{i=1}^k\quot{R}{R\cap P_i^{a_i}}\to \prod_{i=1}^k\quot{R+P_i^{a_i}}{P_i^{a_i}}$ defined by $\phi(r_1+R\cap P_1^{a_1},\dots,r_k+R\cap P_k^{a_k})=(r_1+P_1^{a_1},\dots,r_k+P_k^{a_k})$. We will focus on the first isomorphism.
		
		Assume that $R$ is an ideal-preserving subring of $T$. This immediately tells us that for each $P_i$ dividing $I$, we have $R\cap P_i\nsubseteq P_i^2$. Now let $\phi:\quot{R}{I}\to \prod_{i=1}^k\quot{R}{R\cap P_i^{a_i}}$ be the mapping defined by $\phi(r+I)=(r+R\cap P_1^{a_1},\dots,r+R\cap P_k^{a_k})$. It is plain to see that $\phi$ is a well-defined homomorphism; moreover, $\phi(r+I)=0\implies r\in \bigcap_{i=1}^k R\cap P_i^{a_i}=I$. Then $\ker(\phi)=\{0\}$, so $\phi$ is injective. All that remains is to show that $\phi$ is surjective. To show this, we will show that for each $1\leq i\leq k$, there exists $r_i\in R$ such that $\phi(r_i+I)$ is 0 in every coordinate except the $i^{th}$, in which it is $1+P_i^{a_i}$; that is, $r_i\equiv 1\modulo{P_i^{a_i}}$ and $r_i\in P_j^{a_j}$ for each $j\neq i$.
		
		Since $R$ is an ideal-preserving subring of $T$, we can select some $x_i\in R\cap \prod_{j\neq i}P_j^{a_j}\backslash P_i$. Now since $x_i\notin P_i$, we know that $x_i+I\in U(\sfrac{T}{P_i^{a_i}})$. Let $S\subseteq \quot{T}{P_i^{a_i}}$ be the subset whose cosets contain elements of $R$. Since $R$ is a subring of $T$ and $T$ is integral over $R$, it is plain to see that $S$ is a subring of $\quot{T}{P_i^{a_i}}$ and $\quot{T}{P_i^{a_i}}$ is integral over $S$. Then any element of $U(\sfrac{T}{P_i^{a_i}})$ which lies in $S$ has its inverse lying in $S$; in particular, this means that there exists $y_i\in R$ such that $x_iy_i+P_i^{a_i}=1+P_i^{a_i}$. Then $r_i:=x_iy_i\equiv 1\modulo{P_i^{a_i}}$ and $r_i\in P_j^{a_j}$ for every $j\neq i$. Then $\phi$ is onto and is thus an isomorphism. Therefore, $1\implies 5$.
		
		All that remains is to show that $5\implies 1$; to do so, we will show that $5\implies 4$. Since 5 makes the assumption that for any $P_i|I$, $R\cap P_i\nsubseteq P_i^2$, we only need to show that for any prime $T$-ideals $P_i,P_j$ dividing $I$, $R\cap P_i\nsubseteq P_j$. Using the isomorphism $\phi$ developed above, there must be some $r\in R$ such that $r\in P_i^{a_i}$ and $r\equiv 1\modulo{P_j^{a_j}}$. Then $r\in R\cap P_i\backslash P_j$, so $R\cap P_i\nsubseteq P_j$. Thus, $5\implies 4$, completing the proof. 
	\end{proof}

\end{document}
